Reasoning models can self-jailbreak: their chain of thought transitions from safe to harmful content without adversarial input. Detecting when this shift occurs (the commitment point) requires real-time monitoring, and hidden-state probes offer a direct window into the model's internal reasoning. Yet which transformer layer to read is routinely treated as an implementation detail. We show this choice matters. Across four models (two architectures, two scales), crossing-rate metrics systematically select shallow layers (L0--L10) while threshold-based metrics select deep layers (L14--L63), a divergence of 14 to 61 layers. This gap reverses whether hidden states outperform text: at deep layers, hidden-state probes gain +20--22 percentage points (pp) precision over text (cross-validated at 7--8B); at shallow layers, the two modalities perform equivalently. The mechanism is cosine compression: shallow representations compress directional information, inflating frequency-based metrics near the noise floor. This bias generalizes across models and a second dataset (AdvBench), and the geometric precondition extends to a fifth architecture (Gemma-2). The advantage diminishes with stronger text encoders and remains suggestive at 32B; seven alternative detection methods fail to match supervised probing. Layer selection metric is a critical, overlooked specification choice in safety probing studies.\footnote{Code and data: \url{https://anonymous.4open.science/r/contrastive_trace_safety-A178/}.}
